\documentclass[a4paper,tate,twocolumn,12pt]{ltjtarticle}
\usepackage{luatexja}
\usepackage{luatexja-ruby}
\usepackage[top=1cm,bottom=1cm,left=1cm,right=1cm,landscape]{geometry}

\begin{document}

\thispagestyle{empty}

% Title
\begin{center}
{\LARGE チューリングラブ feat. Sou}

{\large ナナヲアカリ}
\end{center}

\vspace{1em}

\begin{flushleft}

\noindent\begin{minipage}{\linewidth}
あー、\ruby{恋}{こい}の\ruby{定義}{ていぎ}がわかんない

まずスキって\ruby{基準}{きじゅん}もわかんない

\ruby{要}{よう}は、\ruby{恋}{こい}してるときが\ruby{恋}{こい}らしい

\ruby{客観}{きゃっかん}? \ruby{主観}{しゅかん}? エビデンスプリーズ!

\ruby{愛}{あい}は\ruby{計算}{けいさん}じゃ\ruby{解}{と}けない

まず\ruby{普通}{ふつう}の\ruby{計算}{けいさん}も\ruby{解}{と}けない

\ruby{要}{よう}は、そんな\ruby{状態}{じょうたい}が\ruby{愛}{あい}らしい

アイノウ? ユーノウ?
\end{minipage}

\vspace{2em}

\noindent\begin{minipage}{\linewidth}
もう\ruby{大抵}{たいてい}の\ruby{事象}{じしょう}において Q があって A を\ruby{出}{だ}して\ruby{解}{と}けるのに

\ruby{勘違}{かんちがい}って\ruby{間違}{まちが}って \ruby{解}{かい}のないこの\ruby{気持}{きも}ちはなんだろう

(\ruby{検証}{けんしょう} is \ruby{不明瞭}{ふめいりょう})

DAZING!! モーションは\ruby{相対性}{そうたいせい}にステイ
\end{minipage}

\vspace{2em}

\noindent\begin{minipage}{\linewidth}
チューリングラブ \ruby{見}{み}つめあったったって\ruby{解}{と}けないメロウ

ワットイズラブ いま 123 で\ruby{証}{あか}そうか

\ruby{言葉}{ことば}で\ruby{生}{う}み\ruby{出}{だ}すクエスチョン(クエスチョン!)

ハートで\ruby{高鳴}{たかな}るアンサー(アンサー!)

\ruby{測}{はか}ったって\ruby{不確定性}{ふかくていせい} ぼくらの BPM

チューリングラブ \ruby{宙}{ちゅう}に\ruby{舞}{ま}ったって\ruby{信}{しん}じないけど

フォーリンラブ いまは ABC すらバグりそうだ

なんでか\ruby{教}{おし}えてオイラー(オイラー!)

\ruby{感想}{かんそう}きかせてフェルマー(フェルマー!)

\ruby{予測}{よそく}\ruby{不可能}{ふかのう}\ruby{性}{せい} いま\ruby{触}{ふ}れてみる
\end{minipage}

\vspace{2em}

\noindent\begin{minipage}{\linewidth}
さあ\ruby{証明}{しょうめい}しよう(\ruby{証明}{しょうめい}しよう)\ruby{間違}{まちが}いのないように

\ruby{証明}{しょうめい}しよう(\ruby{証明}{しょうめい}しよう)この\ruby{感情}{かんじょう}\ruby{全部}{ぜんぶ}\ruby{全部}{ぜんぶ}

\ruby{証明}{しょうめい}しよう(\ruby{証明}{しょうめい}しよう)\ruby{正解}{せいかい}があるなら

\ruby{証明}{しょうめい}しよう(\ruby{証明}{しょうめい}しよう)シンプルな Q.E.D.

\ruby{今日}{きょう}こそ\ruby{証明}{しょうめい}したい!(しよう!)

\ruby{目}{め}に\ruby{視}{み}えない \ruby{無理難題}{むりなんだい}を ASAP!

\ruby{導}{みちび}き\ruby{出}{だ}したい たどり\ruby{着}{つ}きたい \ruby{最適解}{さいてきかい}へ

\ruby{何度}{なんど}\ruby{空回}{からまわ}ったって ぶつかっちゃったって

\ruby{不可思議}{ふかしぎ}なこのアノマリーを

\ruby{解}{と}き\ruby{明}{あ}かす\ruby{過程}{かてい}さえ\ruby{必要条件}{ひつようじょうけん}です。
\end{minipage}

\vspace{2em}

\noindent\begin{minipage}{\linewidth}
「3\ruby{次元}{じげん}\ruby{(立体)}{かっこりったい}における 2\ruby{点}{てん} AB \ruby{間}{かん}の\ruby{距離}{きょり}を\ruby{求}{もと}めよ。」

\ruby{的}{てき}な\ruby{感}{かん}じだと\ruby{思}{おも}ってたのに

\ruby{座標}{ざひょう}も\ruby{公式}{こうしき}も\ruby{見当}{みあ}たらないし

\ruby{参考}{さんこう}\ruby{資料}{しりょう}にも\ruby{載}{の}っていないから

\ruby{君}{きみ}と\ruby{生}{う}み\ruby{出}{だ}してみせたいじゃん

3\ruby{次元}{じげん}\ruby{(世界)}{かっこせかい}におけるこの\ruby{気持}{きも}ちの\ruby{求}{もと}め\ruby{方}{かた}
\end{minipage}

\vspace{2em}

\noindent\begin{minipage}{\linewidth}
\ruby{○}{まる}だって×だって \ruby{不安}{ふあん}にとってイコールで

\ruby{近}{ちか}くなって\ruby{遠}{とお}くなって \ruby{揺}{ゆ}れるこの\ruby{気持}{きも}ちの\ruby{針}{はり}も(\ruby{条件}{じょうけん} is \ruby{明}{めい}\ruby{瞭}{りょう})

Steppin'!! エモーションは\ruby{相愛}{そうあい}\ruby{性}{せい}ミステイク

チューリングラブ \ruby{騒}{さわ}ぎ\ruby{出}{だ}した\ruby{胸}{むね}の\ruby{痛}{いた}くなるほど

\ruby{証明}{しょうめい}はいまも \ruby{確度}{かくど}を\ruby{増}{ま}しているようだ

\ruby{曖昧}{あいまい}も\ruby{除外}{じょがい}して \ruby{最大}{さいだい}の\ruby{仮説}{かせつ}を

いま、ふたりなら\ruby{実証}{じっしょう}できそう
\end{minipage}

\vspace{2em}

\noindent\begin{minipage}{\linewidth}
チューリングラブ \ruby{見}{み}つめあったったってキリないメロウ

ワットイズラブ いま 123 で\ruby{証}{あか}そうか

\ruby{言葉}{ことば}で\ruby{生}{う}み\ruby{出}{だ}すクエスチョン(クエスチョン!)

ハートで\ruby{高鳴}{たかな}るアンサー(アンサー!)

\ruby{測}{はか}ったって\ruby{不確定性}{ふかくていせい} ぼくらの BPM

チューリングラブ \ruby{宙}{ちゅう}に\ruby{舞}{ま}ったってその\ruby{因果}{いんが}も

フォーリンラブ いまは XYZ まで\ruby{解}{わか}りそうだ

\ruby{確}{たし}かめさせてピタゴラス(ピタゴラス!)

\ruby{確}{たし}かにさせてリーマン(リーマン!)

\ruby{予測}{よそく}\ruby{済}{ず}み\ruby{可変}{かへん}\ruby{性}{せい} いま\ruby{触}{ふ}れてみる
\end{minipage}

\vspace{2em}

\noindent\begin{minipage}{\linewidth}
さあ\ruby{証明}{しょうめい}しよう(\ruby{証明}{しょうめい}しよう)\ruby{間違}{まちが}いのないように

\ruby{証明}{しょうめい}しよう(\ruby{証明}{しょうめい}しよう)この\ruby{感情}{かんじょう}\ruby{全部}{ぜんぶ}\ruby{全部}{ぜんぶ}

\ruby{証明}{しょうめい}しよう(\ruby{証明}{しょうめい}しよう)\ruby{正解}{せいかい}があるなら

\ruby{証明}{しょうめい}しよう(\ruby{証明}{しょうめい}しよう)\ruby{推測}{すいそく}を\ruby{超}{こ}えて!

\ruby{証明}{しょうめい}しよう(\ruby{証明}{しょうめい}しよう)\ruby{間違}{まちが}いのないように

\ruby{証明}{しょうめい}しよう(\ruby{証明}{しょうめい}しよう)この\ruby{感情}{かんじょう}\ruby{全部}{ぜんぶ}\ruby{全部}{ぜんぶ}

\ruby{証明}{しょうめい}しよう(\ruby{証明}{しょうめい}しよう)\ruby{正解}{せいかい}があるなら

\ruby{証明}{しょうめい}しよう(\ruby{証明}{しょうめい}しよう)シンプルな Q.E.D.
\end{minipage}

\end{flushleft}

\end{document}
